\documentclass[reqno]{amsart}
\usepackage[T1]{fontenc}
\usepackage{amsmath,amssymb,amsthm}
\usepackage{xr-hyper}
\usepackage[hidelinks]{hyperref}
\externaldocument[tn-]{20261125T002100--definite-theta-lifts-and-norms__workbook}
\externaldocument[ob-]{20261126T000000--the-objects-of-the-construction__workbook}
\externaldocument[sb-]{20261125T020200--slopes-frames-and-bessel-poincare-series__workbook}
\externaldocument[bp-]{20261126T020210--bessel-poincare-series-and-their-gram-operators__workbook}
\externaldocument[sx-]{20261126T020220--the-seesaw-as-transport-of-relations__workbook}

\newcounter{chap}
% Chapter numbers, generated from 20261125T012000--grh-workbook__guide.org
% by scripts/make-chapters.py.  Do not edit: change the order of the
% index and rerun the script (./preview-site.sh and CI do this).
%
% Usage in a sheet, after \newcounter{chap}:
%   % Chapter numbers, generated from 20261125T012000--grh-workbook__guide.org
% by scripts/make-chapters.py.  Do not edit: change the order of the
% index and rerun the script (./preview-site.sh and CI do this).
%
% Usage in a sheet, after \newcounter{chap}:
%   % Chapter numbers, generated from 20261125T012000--grh-workbook__guide.org
% by scripts/make-chapters.py.  Do not edit: change the order of the
% index and rerun the script (./preview-site.sh and CI do this).
%
% Usage in a sheet, after \newcounter{chap}:
%   \input{chapters}\chapnumber{<sheet ID>}
\makeatletter
\providecommand\chapters@set[2]{\@namedef{chapters@#1}{#2}}
\chapters@set{20261125T001000}{1}% zeros-to-ramp-sums
\chapters@set{20261125T001100}{2}% inert-primes-and-estimate-b
\chapters@set{20261125T005000}{3}% galois-fourier-and-averaging
\chapters@set{20261125T005500}{4}% columns-from-a-sphere
\chapters@set{20261125T000000}{5}% slope-method
\chapters@set{20261126T030100}{6}% a-twist-for-each-prime
\chapters@set{20261125T003000}{7}% the-chain-and-the-square-root
\chapters@set{20261125T010000}{8}% rank-lemma-and-zero-estimate
\chapters@set{20261125T010500}{9}% poincare-series-and-near-orthogonality
\chapters@set{20261124T223000}{10}% elliptic-curves-for-the-proof
\chapters@set{20261125T009500}{11}% fock-model-and-theta-intertwiners
\chapters@set{20261124T223100}{12}% complex-multiplication-and-inert-reduction
\chapters@set{20261125T008000}{13}% one-cm-determinant
\chapters@set{20261124T214500}{14}% supersingular-theta-divisibility
\chapters@set{20261124T214600}{15}% the-isotropic-plane
\chapters@set{20261125T013700}{16}% pairings-on-a-formal-group-of-height-two
\chapters@set{20261125T013500}{17}% hecke-nondegeneracy-on-the-tree
\chapters@set{20261126T030200}{18}% the-bound-is-attained
\chapters@set{20261125T009000}{19}% kronecker-limit-and-chowla-selberg
\chapters@set{20261125T002000}{20}% definite-theta-lifts
\chapters@set{20261125T015500}{21}% siegel-weil-in-the-smallest-case
\chapters@set{20261125T002100}{22}% definite-theta-lifts-and-norms
\chapters@set{20261125T016500}{23}% hecke-operators-and-theta-lifts
\chapters@set{20261125T016700}{24}% divisibility-along-the-chain
\chapters@set{20261125T004000}{25}% denominators-and-sha
\chapters@set{20261125T004100}{26}% tate-twists-and-sha
\chapters@set{20261126T020430}{27}% the-minus-part-by-gauss-sums
\chapters@set{20261125T014100}{28}% congruence-modules
\chapters@set{20261125T014500}{29}% eisenstein-congruences-and-selmer-classes
\chapters@set{20261125T014700}{30}% pseudocharacters-and-the-length-bound
\chapters@set{20261125T011000}{31}% galois-representations-black-box
\chapters@set{20261126T030300}{32}% the-data
\chapters@set{20261126T020210}{33}% bessel-poincare-series-and-their-gram-operators
\chapters@set{20261126T030400}{34}% the-seesaw-identity
\chapters@set{20261126T030500}{35}% divisibility-of-the-composite-lift
\chapters@set{20261126T030600}{36}% theorem-a-short-forms
\chapters@set{20261126T030700}{37}% the-determinant-of-a-chain
\chapters@set{20261126T030800}{38}% the-cost-of-one-lift
\chapters@set{20261126T030900}{39}% denominators-of-projections-to-lifts
\chapters@set{20261126T031000}{40}% twists-with-few-congruences
\chapters@set{20261126T031100}{41}% theorem-c-and-the-main-inequality
\providecommand\chapnumber[1]{%
  \@ifundefined{chapters@#1}%
    {\PackageError{chapters}{Sheet #1 is not listed in the index}%
      {List it in 20261125T012000--grh-workbook__guide.org\MessageBreak
       and run scripts/make-chapters.py.}}%
    {\edef\thechap{\@nameuse{chapters@#1}}}}
\makeatother
\chapnumber{<sheet ID>}
\makeatletter
\providecommand\chapters@set[2]{\@namedef{chapters@#1}{#2}}
\chapters@set{20261125T001000}{1}% zeros-to-ramp-sums
\chapters@set{20261125T001100}{2}% inert-primes-and-estimate-b
\chapters@set{20261125T005000}{3}% galois-fourier-and-averaging
\chapters@set{20261125T005500}{4}% columns-from-a-sphere
\chapters@set{20261125T000000}{5}% slope-method
\chapters@set{20261126T030100}{6}% a-twist-for-each-prime
\chapters@set{20261125T003000}{7}% the-chain-and-the-square-root
\chapters@set{20261125T010000}{8}% rank-lemma-and-zero-estimate
\chapters@set{20261125T010500}{9}% poincare-series-and-near-orthogonality
\chapters@set{20261124T223000}{10}% elliptic-curves-for-the-proof
\chapters@set{20261125T009500}{11}% fock-model-and-theta-intertwiners
\chapters@set{20261124T223100}{12}% complex-multiplication-and-inert-reduction
\chapters@set{20261125T008000}{13}% one-cm-determinant
\chapters@set{20261124T214500}{14}% supersingular-theta-divisibility
\chapters@set{20261124T214600}{15}% the-isotropic-plane
\chapters@set{20261125T013700}{16}% pairings-on-a-formal-group-of-height-two
\chapters@set{20261125T013500}{17}% hecke-nondegeneracy-on-the-tree
\chapters@set{20261126T030200}{18}% the-bound-is-attained
\chapters@set{20261125T009000}{19}% kronecker-limit-and-chowla-selberg
\chapters@set{20261125T002000}{20}% definite-theta-lifts
\chapters@set{20261125T015500}{21}% siegel-weil-in-the-smallest-case
\chapters@set{20261125T002100}{22}% definite-theta-lifts-and-norms
\chapters@set{20261125T016500}{23}% hecke-operators-and-theta-lifts
\chapters@set{20261125T016700}{24}% divisibility-along-the-chain
\chapters@set{20261125T004000}{25}% denominators-and-sha
\chapters@set{20261125T004100}{26}% tate-twists-and-sha
\chapters@set{20261126T020430}{27}% the-minus-part-by-gauss-sums
\chapters@set{20261125T014100}{28}% congruence-modules
\chapters@set{20261125T014500}{29}% eisenstein-congruences-and-selmer-classes
\chapters@set{20261125T014700}{30}% pseudocharacters-and-the-length-bound
\chapters@set{20261125T011000}{31}% galois-representations-black-box
\chapters@set{20261126T030300}{32}% the-data
\chapters@set{20261126T020210}{33}% bessel-poincare-series-and-their-gram-operators
\chapters@set{20261126T030400}{34}% the-seesaw-identity
\chapters@set{20261126T030500}{35}% divisibility-of-the-composite-lift
\chapters@set{20261126T030600}{36}% theorem-a-short-forms
\chapters@set{20261126T030700}{37}% the-determinant-of-a-chain
\chapters@set{20261126T030800}{38}% the-cost-of-one-lift
\chapters@set{20261126T030900}{39}% denominators-of-projections-to-lifts
\chapters@set{20261126T031000}{40}% twists-with-few-congruences
\chapters@set{20261126T031100}{41}% theorem-c-and-the-main-inequality
\providecommand\chapnumber[1]{%
  \@ifundefined{chapters@#1}%
    {\PackageError{chapters}{Sheet #1 is not listed in the index}%
      {List it in 20261125T012000--grh-workbook__guide.org\MessageBreak
       and run scripts/make-chapters.py.}}%
    {\edef\thechap{\@nameuse{chapters@#1}}}}
\makeatother
\chapnumber{<sheet ID>}
\makeatletter
\providecommand\chapters@set[2]{\@namedef{chapters@#1}{#2}}
\chapters@set{20261125T001000}{1}% zeros-to-ramp-sums
\chapters@set{20261125T001100}{2}% inert-primes-and-estimate-b
\chapters@set{20261125T005000}{3}% galois-fourier-and-averaging
\chapters@set{20261125T005500}{4}% columns-from-a-sphere
\chapters@set{20261125T000000}{5}% slope-method
\chapters@set{20261126T030100}{6}% a-twist-for-each-prime
\chapters@set{20261125T003000}{7}% the-chain-and-the-square-root
\chapters@set{20261125T010000}{8}% rank-lemma-and-zero-estimate
\chapters@set{20261125T010500}{9}% poincare-series-and-near-orthogonality
\chapters@set{20261124T223000}{10}% elliptic-curves-for-the-proof
\chapters@set{20261125T009500}{11}% fock-model-and-theta-intertwiners
\chapters@set{20261124T223100}{12}% complex-multiplication-and-inert-reduction
\chapters@set{20261125T008000}{13}% one-cm-determinant
\chapters@set{20261124T214500}{14}% supersingular-theta-divisibility
\chapters@set{20261124T214600}{15}% the-isotropic-plane
\chapters@set{20261125T013700}{16}% pairings-on-a-formal-group-of-height-two
\chapters@set{20261125T013500}{17}% hecke-nondegeneracy-on-the-tree
\chapters@set{20261126T030200}{18}% the-bound-is-attained
\chapters@set{20261125T009000}{19}% kronecker-limit-and-chowla-selberg
\chapters@set{20261125T002000}{20}% definite-theta-lifts
\chapters@set{20261125T015500}{21}% siegel-weil-in-the-smallest-case
\chapters@set{20261125T002100}{22}% definite-theta-lifts-and-norms
\chapters@set{20261125T016500}{23}% hecke-operators-and-theta-lifts
\chapters@set{20261125T016700}{24}% divisibility-along-the-chain
\chapters@set{20261125T004000}{25}% denominators-and-sha
\chapters@set{20261125T004100}{26}% tate-twists-and-sha
\chapters@set{20261126T020430}{27}% the-minus-part-by-gauss-sums
\chapters@set{20261125T014100}{28}% congruence-modules
\chapters@set{20261125T014500}{29}% eisenstein-congruences-and-selmer-classes
\chapters@set{20261125T014700}{30}% pseudocharacters-and-the-length-bound
\chapters@set{20261125T011000}{31}% galois-representations-black-box
\chapters@set{20261126T030300}{32}% the-data
\chapters@set{20261126T020210}{33}% bessel-poincare-series-and-their-gram-operators
\chapters@set{20261126T030400}{34}% the-seesaw-identity
\chapters@set{20261126T030500}{35}% divisibility-of-the-composite-lift
\chapters@set{20261126T030600}{36}% theorem-a-short-forms
\chapters@set{20261126T030700}{37}% the-determinant-of-a-chain
\chapters@set{20261126T030800}{38}% the-cost-of-one-lift
\chapters@set{20261126T030900}{39}% denominators-of-projections-to-lifts
\chapters@set{20261126T031000}{40}% twists-with-few-congruences
\chapters@set{20261126T031100}{41}% theorem-c-and-the-main-inequality
\providecommand\chapnumber[1]{%
  \@ifundefined{chapters@#1}%
    {\PackageError{chapters}{Sheet #1 is not listed in the index}%
      {List it in 20261125T012000--grh-workbook__guide.org\MessageBreak
       and run scripts/make-chapters.py.}}%
    {\edef\thechap{\@nameuse{chapters@#1}}}}
\makeatother

\ifcsname chapters@20261126T020230\endcsname
  \chapnumber{20261126T020230}
\else
  \renewcommand{\thechap}{?}
\fi
\renewcommand{\thesection}{\thechap.\arabic{section}}
\theoremstyle{definition}
\newtheorem{problem}{Problem}[section]
\theoremstyle{plain}
\newtheorem{theorem}{Theorem}[section]
\theoremstyle{remark}
\newtheorem*{leanremark}{In the Lean formalization}
\newcommand{\Z}{\mathbb Z}
\newcommand{\C}{\mathbb C}
\newcommand{\Hc}{\mathcal H}
\newcommand{\Ac}{\mathcal A}
\newcommand{\Ec}{\mathcal E}
\newcommand{\HS}{\mathrm{HS}}
\newcommand{\Bper}{\operatorname{Bper}}
\newcommand{\Pt}{\operatorname{Pt}}
\newcommand{\tl}[1]{\texttt{#1}}
\setlength{\emergencystretch}{3em}

\title{The Bessel identity}
\author{}
\date{November 26, 2026}

\begin{document}
\maketitle
\refstepcounter{chap}\label{chap}

\paragraph{Prerequisites.}
From Chapter~\ref{ob-chap}: the lattices and the period maps \(B_u\) of Section~\ref{ob-sec:lattices}, and the coefficient field \(L\) with its fixed embedding \(\sigma\colon L\to\C\).  From Chapter~\ref{sb-chap}: Bessel realizations~\eqref{sb-eq:bessel}, Theorems~\ref{sb-thm:frame}, \ref{sb-thm:gram} and~\ref{sb-thm:A3}, and Problems~\ref{sb-pr:ramified} and~\ref{sb-pr:kernel}.  From Chapter~\ref{bp-chap}: the Poincar\'e series \(P_u\), the spaces \(C_u\), the toy of Section~\ref{bp-sec:smallest} and Theorem~\ref{bp-thm:near}.  From Chapter~\ref{sx-chap}: Problems~\ref{sx-pr:anchor}, \ref{sx-pr:compose}, \ref{sx-pr:unit-ratio} and~\ref{sx-pr:toy-chain}, and Theorem~\ref{sx-thm:bottom}.  From Chapter~\ref{tn-chap}: the period of Section~\ref{tn-sec:seesaw}.  Inner products on complex Hilbert spaces are conjugate-linear in the first variable, as in Chapter~\ref{sb-chap}.

\medskip
Condition~(d) of Theorem~\ref{sb-thm:A3} asks that the period maps have a Bessel realization: each \(B_u\) is, up to unitary changes of variable and a scalar, an inner product with a Poincar\'e series.  Chapter~\ref{sx-chap} moved the pairing at the top of the chain to a Bessel period at the bottom, and Chapter~\ref{bp-chap} wrote Bessel periods as inner products with Poincar\'e series.  This chapter assembles these into the identity that the Lean formalization \cite{Lean} proves, and says what it gives.

\section{The statement}\label{sec:statement}

For one record, let \(\Hc\) be the space of automorphic forms of Chapter~\ref{bp-chap}, of dimension \(D\).  Let \(\Ec_0:=(R\otimes L)\otimes_{L,\sigma}\C\) be the complexification of the row space of Chapter~\ref{ob-chap}, of dimension \(D_{\rm r}\le D\), with the inner product given by the norm of \(R\) at \(\sigma\), and let \(s\colon\Ec_0\to\Hc\) be a conjugate-linear isometry, the row embedding.  For each word \(u\) let \(M_u\subset C_u\) be the complexification of the forms of level \(0\) of \(u\) ramified at \(T_u\), with the inner product of \(C_u\), and \(P_u\colon C_u\to\Hc\) the Bessel--Poincar\'e map.  The period map gives a \(\C\)-bilinear pairing \(B_u\colon M_u\times\Ec_0\to\C\); on vectors coming from the \(L\)-structures it is defined, as in (B) of Section~\ref{ob-sec:lattices}, by the pairing at the end of the chain, and \(B_u^{\rm alg}\) denotes the corresponding \(L\)-valued pairing of the lattices.

\begin{theorem}[Bessel identity]\label{thm:bessel}
For every record of the family at a large scale:
\begin{enumerate}
\item[(i)] \(B_u(\varphi,F)=\sigma\bigl(B_u^{\rm alg}(\varphi,F)\bigr)\) for \(\varphi\) and \(F\) in the \(L\)-structures;
\item[(ii)] there are a unitary operator \(V\) of \(\Ec_0\) and, for each word \(u\), a number \(a_u\ne0\) and a unitary operator \(U_u\) of \(M_u\) with
\[
 B_u(\varphi,F)=a_u\,\langle sVF,\ P_uU_u\varphi\rangle_\Hc\qquad(\varphi\in M_u,\ F\in\Ec_0).
\]
\end{enumerate}
\end{theorem}
\noindent\emph{Status}: proved in the formalization, for the family it constructs; the proof is the subject of this chapter, and we say at each step how far we have followed it.  In \cite{Adr} the statement is Theorem~\ref{sb-thm:bessel} of Chapter~\ref{sb-chap}, derived there from hypotheses assumed there, among them a seesaw relation for an abstract period map, which the formalization does not state.

\begin{problem}\label{pr:conj}
The right side of~(ii) is \(\C\)-linear in \(\varphi\) and in \(F\).  If \(s\) were \(\C\)-linear and~(ii) held with \(B_u\) bilinear, then \(B_u=0\).
\end{problem}

\begin{problem}\label{pr:realization}
By Theorem~\ref{thm:bessel}(i), \(B_u\) is the \(\C\)-bilinear extension of \(\sigma\circ B_u^{\rm alg}\), that is, the map \(B_{u,v}\) of Chapter~\ref{sb-chap} at the place \(v\) of \(\sigma\); so Theorem~\ref{thm:bessel}(i) and~(ii) give a Bessel realization~\eqref{sb-eq:bessel} of the maps \(B_u\), with \(Y_v=\Ec_0\), \(V_{u,v}=M_u\), \(\iota_0=s\), \(U=V\), \(\alpha_u=a_u\), \(\iota_u\) the inclusion \(M_u\subset C_u\), and the same \(U_u\).
\end{problem}

In the toy of Section~\ref{bp-sec:smallest}, let \(\Ec_0:=\Ac_{(a)}(2)\), with \(s(F)\) the form whose coefficients in the basis \(e^\nu\) are the complex conjugates of those of \(F\), let \(M_j=C_j:=\Ac_{(j)}(1)\), and put \(B_j(\xi,F):=B(sF,\xi)\), with the period \(B\) of Section~\ref{bp-sec:poincare} read with the inner product of this chapter, so that \(B(\psi,\xi)\) is conjugate-linear in \(\psi\); it is the complex conjugate of the period \(P(\psi,\xi)\) of Chapter~\ref{tn-chap} when the latter is read as linear in \(\psi\).

\begin{problem}\label{pr:toy}
\(s\) is a conjugate-linear isometry of \(\Ac_{(a)}(2)\), \(B_j\) is bilinear, and \(B_j(\xi,F)=\langle sF,P\xi\rangle\): Theorem~\ref{thm:bessel}(ii) holds with \(a_j=1\) and \(U_j=V=1\).  For \(a=8\): \(B_0(\xi_0,e_1^8+e_2^8)=1/4\) and \(B_4(\xi_4,e_1^4e_2^4)=1/280\).
\end{problem}

Problem~\ref{sx-pr:toy-chain} of Chapter~\ref{sx-chap} adds the chain: for \(F=\Theta_{(a)}\varphi\) with \(\varphi\) real, the period at the top equals \(B_j(\xi,F)\).

\section{From the kernel identity to the Bessel identity}\label{sec:assembly}

In the formalization the row space is reached from below: there are a finite-dimensional inner product space \(V_0\) of forms and a linear map \(\iota\colon V_0\to\Ec_0\) with \(\iota(V_0)=\Ec_0\), and \(B_u(\varphi,\iota q)=\mathrm{per}_u(\varphi,q)\) for \(q\in(\ker\iota)^\perp\), where \(\mathrm{per}_u\) is the pairing at the end of the chain of \(u\) at the base twist.  The \emph{row dictionary} is an isometry \(j\) of \(\Ec_0\) into \(L^2\) of the automorphic quotient with \(s(F)=\overline{j(F)}\).  The \emph{column dictionary} is a unitary operator \(W_u\) of \(M_u\): the column of the definite side attached to \(\psi\in M_u\) is the element \(W_u\psi\) of \(M_u\subset C_u\).  In Sections~\ref{sec:bases} and~\ref{sec:mult}, \(P_u\psi\) for \(\psi\in M_u\) means \(P_uW_u\psi\).

\begin{theorem}[kernel identity]\label{thm:kernel}
There are a unitary operator \(V\) of \(\Ec_0\) and, for each word \(u\), a number \(a_u\ne0\) and a unitary operator \(U_u\) of \(M_u\) such that
\[
 \mathrm{per}_u(\varphi,q)=a_u\,\bigl\langle\overline{j(V\iota q)},\ P_uW_uU_u\varphi\bigr\rangle
\]
for all \(\varphi\in M_u\) and \(q\in V_0\).
\end{theorem}
\noindent\emph{Status}: proved in the formalization from Theorem~\ref{thm:mult} below by the argument of Section~\ref{sec:bases}.

\begin{problem}\label{pr:lift-rows}
Let \(\iota\colon V_0\to\Ec_0\) be a linear map of finite-dimensional inner product spaces with \(\iota(V_0)=\Ec_0\).  For every \(x\in\Ec_0\) there is \(q\in(\ker\iota)^\perp\) with \(\iota q=x\).
\end{problem}

\begin{problem}\label{pr:assembly}
Theorem~\ref{thm:kernel}, the two dictionaries and \(B_u(\varphi,\iota q)=\mathrm{per}_u(\varphi,q)\) imply Theorem~\ref{thm:bessel}(ii), with \(U_u\) replaced by \(W_uU_u\).
\end{problem}

\section{Bilinear forms on bases}\label{sec:bases}

The step behind Theorem~\ref{thm:kernel} compares two bilinear forms on coordinate vectors.  Fix an \(L\)-basis \(m_1,\dots,m_d\) of the \(L\)-structure of \(M_u\), and write \(\varphi\in L^d\) also for \(\sum_l\varphi_lm_l\in M_u\).  The pairing \(\mathrm{per}_u\) is \(\C\)-bilinear on \(M_u\times V_0\), and \(V_0\) has an \(L\)-structure with a basis \(r_1,\dots,r_n\), which is a \(\C\)-basis of \(V_0\); for \(F\in L^n\) put \(\mathrm{row}(F):=\sum_kF_kr_k\in V_0\).  For each twist \(t\) there are row coordinates \(F\in L^{n_t}\) (with \(n_{t_0}=n\)), a \(\C\)-linear map \(\Theta_t\) from \(\C^{n_t}\) to \(\Ec_0\), and an \(L\)-bilinear period \(\Pt_{u,t}\colon L^d\times L^{n_t}\to L\) of the definite side; the period map at the twist \(t\) on coordinates is \(B_{u,t}(\varphi,F)=\sigma\Pt_{u,t}(\varphi,F)\), and at the base twist \(\mathrm{per}_u(\varphi,\mathrm{row}(F))=\sigma\Pt_{u,t_0}(\varphi,F)\) (the \emph{chain tie}).  Finally \(\iota(\mathrm{row}(F))=\lambda^{-1}\Theta_{t_0}(F)\) for a number \(\lambda\ne0\) (the \emph{row tie}).

\begin{theorem}[multiplicity identity]\label{thm:mult}
There are a unitary operator \(V\) of \(\Ec_0\) and, for each word \(u\), numbers \(a_u\ne0\), a unitary operator \(U_u\) of \(M_u\), and numbers \(\zeta(u,t)\ne0\) with \(\zeta(u,t_0)=1\), such that
\[
 \zeta(u,t)\,\sigma\bigl(\Pt_{u,t}(\varphi,F)\bigr)=a_u\,\bigl\langle sV\Theta_t(F),\ P_uU_u\varphi\bigr\rangle
\]
for all twists \(t\) and all coordinate vectors \(\varphi\in L^d\), \(F\in L^{n_t}\).  Moreover, if \(\zeta(u,t)=\sigma(a)\) with \(a\in L\), then \(v_w(a)=0\) at every finite place \(w\) of \(L\) above no prime of the set \(S_t\) of Chapter~\ref{sx-chap}.
\end{theorem}

\begin{problem}\label{pr:extend}
Let \(\beta\colon L^d\times L^n\to L\) be \(L\)-bilinear and \(\Phi\colon\C^d\times\C^n\to\C\) be \(\C\)-bilinear, with \(\Phi(e_i,e_k)=\sigma\beta(e_i,e_k)\) on the standard basis vectors.  Then \(\Phi(\sigma x,\sigma y)=\sigma\beta(x,y)\) for all \(x\in L^d\), \(y\in L^n\), and \(\Phi\) is determined by \(\beta\).
\end{problem}

\begin{problem}\label{pr:kernel-from-mult}
Assume Theorem~\ref{thm:mult}, the chain tie and the row tie.  Then the two \(\C\)-bilinear forms \((\varphi,q)\mapsto\mathrm{per}_u(\varphi,q)\) and \((\varphi,q)\mapsto a_u\lambda\langle sV\iota q,P_uU_u\varphi\rangle\) on \(M_u\times V_0\) agree on \((m_l,r_k)\), hence everywhere; this is Theorem~\ref{thm:kernel} with \(a_u\lambda\) in place of \(a_u\).
\end{problem}

In the formalization the embedding in the chain tie may be the complex conjugate of \(\sigma\), and then the scalars are conjugated accordingly; we suppress this.

\section{The multiplicity identity}\label{sec:mult}

Each \(e\in\Ec_0\) has a theta lift \(\mathsf R_e\), an automorphic form on the indefinite group at the top, whose restriction to the base is a row form \(T_e\) in the sense of Chapter~\ref{sx-chap}; every row form at the base is some \(T_e\), and \(\mathrm{row}_t(F)\) of Chapter~\ref{sx-chap} is \(T_{\Theta_t(F)}\).  Recall that \(P_ux\) means \(P_uW_ux\).  Theorem~\ref{thm:mult} comes from five inputs.
\begin{itemize}
\item[(E)] There are a vector \(v\) of the archimedean discrete series and an isometric bijection \(J\) from the space \(\{\mathsf R_e:e\in\Ec_0\}\) onto \(\Hc\), such that \(J(\mathsf R_e)\) is almost everywhere the complex conjugate of the evaluation of \(\mathsf R_e\) at \(v\).  \emph{Status}: proved in the formalization from a proposition stating a classification of discrete series of unitary groups \cite{HC}; we have not traced how that proposition is supplied.
\item[(Pa)] There are a unitary operator \(V\) of \(\Ec_0\), induced by translation by a central element, and \(\kappa\ne0\) such that \(\langle sVe,P_ux\rangle=\kappa\,b(e,x)\) for \(e\in\Ec_0\) and \(x\in M_u\), where \(b(e,x):=\langle J(\mathsf R_e),P_ux\rangle_\Hc\) is the integral of the evaluated theta lift of \(e\) against the Poincar\'e series of \(x\), without complex conjugation.  \emph{Status}: proved in the formalization; we have read only its statement.
\item[(J)] Theorem~\ref{sx-thm:bottom} of Chapter~\ref{sx-chap}, with its unitary operator \(U\) of \(M_u\), its constants \(c_t\) and, for each row form \(T\), its function \(p\).
\item[(Bo)] There is a unitary operator \(R_u\) of \(M_u\) such that for all \(e,e'\in\Ec_0\), \(x,x'\in M_u\) and \(\alpha,\alpha'\in\C\): if \(\alpha\Bper(x,T_e)+\alpha'\Bper(x',T_{e'})=0\), then \(\alpha b(e,R_ux)+\alpha'b(e',R_ux')=0\).  \emph{Status}: proved in the formalization; by our reading it is the unfolding of Chapter~\ref{bp-chap}, but we have not read its proof.
\item[(NZ)] For every word \(u\) there is \(\varphi\in M_u\) with \(P_u\varphi\ne0\).  \emph{Status}: an assumption of the step of the formalization that proves Theorem~\ref{thm:mult}; we have not traced where it is proved.
\end{itemize}

\begin{problem}\label{pr:nz}
Assume the second inequality of Theorem~\ref{bp-thm:near} with \(e_0\le\frac12\), and \(\dim M_u>\frac12\dim C_u\).  Then (NZ) holds.
\end{problem}

\begin{problem}\label{pr:anchor-pair}
(E) and (NZ) give, for every word \(u\), vectors \(e^\circ\in\Ec_0\) and \(x^\circ\in M_u\) with \(b(e^\circ,x^\circ)\ne0\).
\end{problem}

\begin{problem}\label{pr:prop}
Assume (J), (Bo) and the conclusion of Problem~\ref{pr:anchor-pair}, and put \(b(t,\varphi,F):=b(\Theta_t(F),R_uU\varphi)\).  There are numbers \(p_0\ne0\) and \(b_0\ne0\), depending on \(u\), such that \(c_t\,B_{u,t}(\varphi,F)\,b_0=p_0\,b(t,\varphi,F)\) for all \(t\), \(\varphi\), \(F\).
\end{problem}

\begin{problem}\label{pr:mult-assembly}
Assume Problem~\ref{pr:prop} and (Pa).  Then the identity of Theorem~\ref{thm:mult} holds with \(\zeta(u,t)=c_t/c_{t_0}\), \(a_u=p_0/(\kappa b_0c_{t_0})\) and \(U_u:=R_uU\).
\end{problem}

\begin{problem}\label{pr:zeta}
In Problem~\ref{pr:mult-assembly}, \(\zeta(u,t_0)=1\) and \(\zeta(u,t)\ne0\); and if \(\zeta(u,t)=\sigma(a)\) with \(a\in L\), then \(v_w(a)=0\) at every finite place \(w\) of \(L\) above no prime of the set \(S_t\) of Chapter~\ref{sx-chap}.
\end{problem}

Problems~\ref{pr:anchor-pair}--\ref{pr:zeta} prove Theorem~\ref{thm:mult} from (E), (Pa), (J), (Bo) and (NZ); Problems~\ref{pr:extend} and~\ref{pr:kernel-from-mult} give Theorem~\ref{thm:kernel}; and Problems~\ref{pr:lift-rows} and~\ref{pr:assembly} give Theorem~\ref{thm:bessel}(ii).  Part~(i) is proved in the formalization from the same multiplicity identity and the chain tie, by an argument we have not followed.

\section{What the identity gives}\label{sec:use}

Fix orthonormal bases \((b_k)\) of \(\Ec_0\), \((e_l)\) of \(M_u\) and \((f_i)\) of \(\Hc\).  In the bases of the \(L\)-structures the matrix of \(B_u\) differs from the one below by invertible matrices, the \(Y_i\) and \(Y'\) of Theorem~\ref{sb-thm:frame}.

\begin{problem}\label{pr:matrix}
Let \(\mathsf S_{ki}:=\langle sVb_k,f_i\rangle\) and \(\Pi_{il}:=\langle f_i,P_uU_ue_l\rangle\).  Then \(B_u(e_l,b_k)=a_u(\mathsf S\Pi)_{kl}\), and \(\mathsf S\) is a \(D_{\rm r}\times D\) matrix with orthonormal rows.
\end{problem}

Complete \(\mathsf S\) to a unitary \(D\times D\) matrix \(\hat{\mathsf S}\), and put \(A_u:=\hat{\mathsf S}\Pi\), whose first \(D_{\rm r}\) rows are \(a_u^{-1}\) times the matrix of \(B_u\).

\begin{problem}\label{pr:gram}
With \(\gamma_u:=c_0(u)^{-1}\) and \(K_{uu}\) as in Chapter~\ref{bp-chap}, \(\gamma_uA_u^*A_u-1=U_u^*(\iota_u^*K_{uu}\iota_u-1)U_u\), where \(\iota_u\colon M_u\to C_u\) is the inclusion; and \(Q:=\sum_u\gamma_uA_uA_u^*\) is unitarily equivalent to \(\sum_u\gamma_uP_u\iota_u\iota_u^*P_u^*\).  So the hypotheses of Theorem~\ref{sb-thm:gram} for the matrices \(\hat{\mathsf S}P_u\) and \(V_u:=\iota_uU_u\) are those for \(P_u\) and \(\iota_uU_u\).
\end{problem}

\begin{problem}\label{pr:rank}
Assume the second inequality of Theorem~\ref{bp-thm:near}.  Then the rank of \(B_u\) is at least \(\dim M_u-e_0\dim C_u-(D-D_{\rm r})\).
\end{problem}

With Theorem~\ref{bp-thm:near} Problem~\ref{pr:gram} gives conditions~(v) and~(vi) of Theorem~\ref{sb-thm:frame}, through Theorem~\ref{sb-thm:gram}, and Problem~\ref{pr:realization} gives~(d) of Theorem~\ref{sb-thm:A3}.

\begin{leanremark}
\raggedright In \cite{Lean}, Theorem~\ref{thm:bessel} is \tl{BesselId}, a clause of \tl{LinkDIV\_at}; part~(i) is \tl{D7.BuRational}, proved through \tl{buRat\_rrFull\_closed\_alt}, and part~(ii) is \tl{adm13}, with \(\langle sF,P_u\varphi\rangle\) the definition \tl{D7.besS}.  The period at the end of the chain is \tl{perCol} (\tl{PtilCol} of \tl{ColumnChain.col}), tied to \(B_u\) by \tl{Rows.BuDefQ}.  Problems~\ref{pr:lift-rows} and~\ref{pr:assembly} are \tl{exists\_Q\_iota} and \tl{besselId\_of\_joins}, reached from \tl{adm13\_of\_kernel\_join} and \tl{adm13\_of\_colJoin}; Theorem~\ref{thm:kernel} is \tl{KerJoinS}, and Problem~\ref{pr:kernel-from-mult} is \tl{kernelJoin\_of\_multiplicity}.  Theorem~\ref{thm:mult} is \tl{MIstar2} (with \tl{MIid}, \tl{MIUnit2}, \tl{AuNZ}), proved by \tl{miStar2\_at\_realisation\_alt} and \tl{miStar2\_of\_anchored}.  The inputs are \tl{exists\_conjVec} (E), \tl{pairForm\_alt} (Pa), \tl{JointCondition} (J), \tl{botRel\_of\_pin\_of\_mem\_center} (Bo) and \tl{PuNZ} (NZ); Problem~\ref{pr:prop} is \tl{anchor\_pairs\_alt} with \tl{anchor\_word\_alt}.  Problems~\ref{pr:matrix} and~\ref{pr:gram} correspond to \tl{blockReal\_at\_of} and \tl{matK\_at\_of}.
\end{leanremark}

\section*{Hints}

\subsection*{\S\ref{sec:statement}}
\begin{description}
\item[\ref{pr:conj}] Replace \(F\) by \(\mathrm iF\).
\item[\ref{pr:toy}] Problem~\ref{bp-pr:unfold}; Problem~\ref{bp-pr:Pj} and \(\lVert e^\nu\rVert^2=\nu!/a!\).
\end{description}

\subsection*{\S\ref{sec:assembly}}
\begin{description}
\item[\ref{pr:lift-rows}] Subtract from any preimage its projection to \(\ker\iota\).
\item[\ref{pr:assembly}] Write the second argument as \(\iota q\) with \(q\in(\ker\iota)^\perp\), by Problem~\ref{pr:lift-rows}.
\end{description}

\subsection*{\S\ref{sec:bases}}
\begin{description}
\item[\ref{pr:extend}] Expand in the bases.
\item[\ref{pr:kernel-from-mult}] Problem~\ref{pr:extend} at \(t=t_0\), where \(\zeta=1\).
\end{description}

\subsection*{\S\ref{sec:mult}}
\begin{description}
\item[\ref{pr:nz}] Problem~\ref{sb-pr:kernel}.
\item[\ref{pr:anchor-pair}] \(J\) is onto.
\item[\ref{pr:prop}] Take \(T:=T_{e^\circ}\) in (J).
\item[\ref{pr:zeta}] Theorem~\ref{sx-thm:bottom}.
\end{description}

\subsection*{\S\ref{sec:use}}
\begin{description}
\item[\ref{pr:matrix}] Expand \(P_uU_ue_l\) in the basis \((f_i)\); \(s\) is a conjugate-linear isometry.
\item[\ref{pr:gram}] \(\hat{\mathsf S}^*\hat{\mathsf S}=1\).
\item[\ref{pr:rank}] Problems~\ref{pr:matrix} and~\ref{sb-pr:kernel}; \(\mathsf S\) has a kernel of dimension \(D-D_{\rm r}\).
\end{description}

\begin{thebibliography}{Lean}
\bibitem[Adr]{Adr} J.~Adriazola, \emph{The Alp\"oge proof of the Riemann hypothesis}, manuscript, revised 24 November 2026, \url{https://github.com/jimjamzola/RH-Alpoge-Reconstruction}.
\bibitem[HC]{HC} Harish-Chandra, \emph{Discrete series for semisimple Lie groups.~II}, Acta Math.\ \textbf{116} (1966), 1--111.
\bibitem[Lean]{Lean} The Lean community, \emph{riemann-hypothesis}, a Lean formalization of Alp\"oge's argument, \url{https://github.com/leanprover-community/riemann-hypothesis}.
\end{thebibliography}
\end{document}
